\documentclass[12pt, a4paper]{article}

\usepackage{ctex}
\usepackage{cite}
\usepackage{amsmath, amssymb}
\usepackage{graphicx}
\usepackage{hyperref}
\usepackage{booktabs}
\usepackage{geometry}

\geometry{margin=2.5cm}

\hypersetup{
    colorlinks=true,
    linkcolor=blue,
    citecolor=blue,
    urlcolor=blue
}

\begin{document}

\title{\textbf{AI赋能高校地理科学专业GIS实践课的创新案例设计}}

\newpage

\section{摘要}

针对高校地理科学专业GIS实践课"重操作轻思维"、与科研实践脱节的困境，本文设计并论述了基于生成式AI的城市物流园区选址分析教学案例，构建"问题定义—数据获取—空间建模—成果表达"四阶段人机协同教学模式。该模式以AI为技术协作者、教师为教学设计者，形成"学生—AI—教师"互动闭环，重点培养学生AI辅助GIS空间分析能力及科研创新能力。教学评价采用过程性评价（40\%）、成果性评价（40\%）与反思性评价（20\%）相结合的综合模式。创新点包括：从"操作导向"向"思维导向"的范式转变、"双师协同"教学模式重构、人机协同素养培育。本研究为培养适应新时代需求的复合型地理专业人才提供了可借鉴的教学模式。

\textbf{关键词}：生成式AI；GIS实践课；人机协同；地理科学专业；教学创新

\newpage

\section{一、引言与文献综述}

\subsection{（一）研究背景}

在教育数字化转型背景下，高校地理科学专业承担着培养具备GIS与AI素养的复合型专业人才的使命。地理空间人工智能（GeoAI）正深刻改变地理数据的分析与应用方式\citep{Jiang2023GeoAI}，GIS空间分析能力已成为地理专业人才的核心竞争力。

然而，传统GIS实践教学存在明显痛点：教学内容偏重软件操作，学生"会点菜单，不会建模"\citep{杜清运2017}；空间分析过程抽象，缺乏真实项目驱动；教学内容与科研实践脱节\citep{王民2018}。生成式AI的快速发展为解决上述问题提供了新思路——大语言模型可作为"智能协作者"辅助代码生成、模型构建与结果解读\citep{Dwivedi2023}。

\subsection{（二）文献综述}

近年来，AI赋能GIS教育已成为研究热点。\citet{Zhang2023AI}研究表明，AI辅助工具可显著提升GIS学习效率与深度；\citet{Wang2020}指出AI语言模型在GIS教育中具有应用潜力；\citet{Zhao2021}探索了人工智能赋能地理实践教学的模式与路径；\citet{Luo2022}提出了人机协同视域下的地理智慧教育理论框架。在国际领域，GeoAI教育研究聚焦于如何将机器学习与空间分析有机整合\citep{Klosterman2015}，以及GIS空间分析的教学内容与方法改革\citep{杜清运2017}。

已有研究存在以下不足：（1）缺乏针对AI与GIS实践深度融合的可操作性教学设计；（2）评价体系研究薄弱，尤其是AI辅助教学的过程性评价；（3）较少关注AI时代地理人才核心能力的重新定义。本研究针对上述不足，设计了基于生成式AI的城市物流园区选址分析教学案例，构建了四阶段人机协同教学模式。

\subsection{（三）研究贡献}

本文贡献包括：（1）构建"问题定义—数据获取—空间建模—成果表达"四阶段人机协同教学模式；（2）设计针对高校地理专业特点的AI辅助GIS教学流程与Prompt示例；（3）建立突出过程性评价的多元化教学评价体系。

\newpage

\section{二、教学目标与核心素养}

基于建构主义教学理念与认知负荷理论\citep{李家清2015}，本案例设定三维教学目标：

\textbf{（1）知识目标}：掌握GIS空间分析基本原理（缓冲区分析、叠加分析、网络分析、适宜性评价）；理解生成式AI在地理数据处理与空间建模中的应用逻辑；掌握多因子加权评价模型的构建方法。

\textbf{（2）能力目标}：能运用AI工具辅助完成城市选址类项目的完整GIS分析流程；具备利用AI生成空间分析代码、解决实际问题的能力；能对AI生成结果进行批判性评估与修正优化；能够撰写规范的GIS分析报告。

\textbf{（3）素养目标}：培养"人机协同"思维，提升批判性使用AI工具的意识；强化地理实践力与综合思维；形成严谨的科研态度与空间思维能力；为未来从事地理信息相关科研工作奠定基础。

核心素养对接方面，本案例紧扣地理科学专业培养要求，注重培养学生的人地协调观、区域认知能力和地理实践力\citep{教育部2022}。

\newpage

\section{三、教学案例设计}

\subsection{（一）案例背景与任务}

以某地级市（人口约300万）为研究区域，要求学生为该市规划一座"生态友好型城市物流园区"。物流园区选址是典型的多因子适宜性评价问题，涉及以下因素的综合权衡：

\begin{itemize}
\item \textbf{交通条件}：高速公路距离、国省道距离、交通枢纽可达性
\item \textbf{地形地貌}：地面坡度、地表起伏度、地质灾害易发性
\item \textbf{生态环境}：自然保护区缓冲区、重要水体距离、生态红线、湿地保护范围
\item \textbf{人口社会}：居民点距离、耕地占用情况
\item \textbf{土地利用}：地价水平、现状用地类型
\end{itemize}

案例的复杂性在于：各因子之间存在冲突（如交通便利区域往往地价较高），学生需通过科学的权重确定与叠加分析方法在多目标间寻求平衡。

\subsection{（二）技术工具与数据}

\textbf{GIS平台}：ArcGIS Pro或QGIS作为空间分析主平台。

\textbf{Python环境}：Jupyter Notebook结合ArcPy/GeoPandas实现空间数据处理。

\textbf{生成式AI}：ChatGPT、文心一言或通义千问作为"智能协作者"，辅助代码生成与问题讨论。

\textbf{数据源}：OpenStreetMap（道路网、居民点、水系）、地理空间数据云（DEM数据）、教师提供的区域基础地理数据。

\subsection{（三）教学流程框架}

本案例采用"学生—AI—教师"三方互动的"人机协同四阶段"教学模式，各阶段层层递进、环环相扣。

\newpage

\section{四、教学过程设计：人机协同四阶段}

\subsection{阶段一：问题定义与AI辅助方案设计（8课时）}

\textbf{核心任务}：学生分组（4-5人/组），明确选址目标与约束条件，界定"生态友好型"内涵，识别关键利益相关方，初步确定评价因子及优先级。

\textbf{AI介入方式}：AI辅助生成结构化分析框架。

\textbf{Prompt示例}：
\begin{itemize}
\item "请为'生态友好型物流园区'GIS选址分析设计完整技术框架，包括：数据需求清单、分析步骤流程、评价指标体系、权重确定方法、结果验证方法。"
\item "在多因子加权评价中，常用的权重确定方法有哪些？各有什么优缺点？推荐适合本科教学的方法并说明步骤。"
\item "请提供国家层面生态保护区划分标准（自然保护区、生态红线、湿地保护范围、河湖蓝线等），哪些应设为'绝对排除区'？"
\end{itemize}

\textbf{教学要点}：引导学生批判性评估AI方案，结合地域实际调整指标权重。例如西南山区需提高地质灾害敏感性权重，东部平原则需更关注洪水风险。培养学生专业判断能力是本阶段关键。

\textbf{成果}：《物流园区选址分析方案》（含数据需求清单、技术路线、评价指标体系、权重方案）。

\subsection{阶段二：数据获取与AI辅助处理（8课时）}

\textbf{核心任务}：获取OSM道路网/DEM数据，完成投影转换、拓扑检查，将DEM转换为坡度图并按适宜性标准重分类为5级。

\textbf{AI介入方式}：AI生成数据获取与预处理Python代码，学生在Jupyter Notebook运行验证。

\textbf{Prompt示例}：
\begin{itemize}
\item "如何使用Python从OpenStreetMap获取某市道路网数据？请给出使用osmnx库的完整代码示例，包括安装、数据下载、范围指定、等级分类、本地保存。"
\item "请写Python代码实现：读取DEM $\rightarrow$ 计算坡度 $\rightarrow$ 按[0-5°)、[5-15°)、[15-25°)、[25-35°)、[>35°]重分类为5级 $\rightarrow$ 保存结果。使用GeoPandas或Rasterio库。"
\end{itemize}

\textbf{教学要点}：AI代码可能存在边界条件处理不当、投影转换错误等问题，学生需理解代码逻辑才能发现修正。强调数据质量评估与一致性处理的重要性。

\textbf{成果}：预处理后的矢量/栅格数据集 + 数据处理日志（含AI交互记录）。

\subsection{阶段三：空间建模与AI协同优化（12课时）}

\textbf{核心任务}：构建多因子加权评价模型，包括：各因子缓冲区/距离图生成、图层栅格化、AHP权重确定、加权叠加分析、敏感性分析。

\textbf{AI介入方式}：AI生成空间建模代码；提供"反事实分析"支持——快速生成对比方案，帮助理解不同参数设置对结果的影响。

\textbf{Prompt示例}：
\begin{itemize}
\item "请用ArcPy或GeoPandas写多因子加权叠加分析代码：输入5个重分类栅格（值域1-5），权重[0.25, 0.20, 0.20, 0.20, 0.15]，计算加权求和并分5级输出，需含完整注释。"
\item "假设生态保护区因子权重从0.25提高到0.35，预测选址结果如何变化？从地理学角度分析说明什么问题？"
\item "如何用成本距离（Cost Distance）分析替代直线距离，更准确地评估交通可达性？请给出Python代码示例。"
\end{itemize}

\textbf{反事实分析示例}：当生态保护区权重提高后，原30\%中适宜区变为低适宜区，选址重心向东南偏移约2km。这说明交通效率与生态保护存在矛盾，需决策者权衡。反事实分析使学生能在有限课时内进行更多情景模拟，深化对政策敏感性的理解。

\textbf{成果}：选址适宜性评价图（5级）、权重合理性论证、反事实分析报告。

\subsection{阶段四：成果表达与AI辅助科研素养培育（8课时）}

\textbf{核心任务}：完善选址分析报告，撰写规范的GIS专业研究报告，撰写AI辅助GIS实践反思报告。

\textbf{AI介入方式}：AI辅助生成专业报告框架与科研规范指导。

\textbf{Prompt示例}：
\begin{itemize}
\item "请提供一份规范的GIS选址分析报告结构框架，包括：摘要、关键词、引言、研究区概况、数据与方法、结果分析、讨论与结论、参考文献等部分的具体写作要点。"
\item "如何规范地撰写GIS分析报告中的地图制图说明？请提供地图要素（标题、图例、比例尺、指北针、坐标系统）规范表达的示例。"
\item "提供一个撰写《AI与地理实践》反思报告的结构框架，引导学生分析：AI在哪些环节提供有效帮助？产生哪些误导或困扰？人机协同边界在哪？对未来地理研究中AI应用的展望？"
\end{itemize}

\textbf{教学要点}：反思报告是课程关键环节，引导学生从效率、理解、局限、伦理四维度审视AI工具，形成客观全面的认知。重点培养学生严谨的科研态度和批判性思维。

\textbf{成果}：（1）选址分析报告（PDF，含规范地图、权重表、结果分析）；（2）《AI辅助GIS实践反思报告》（Word）。

\newpage

\section{五、教学评价体系}

本案例采用过程性评价（40\%）、成果性评价（40\%）与反思性评价（20\%）相结合的多元化评价体系，如表1所示。

\begin{table}[htbp]
\centering
\caption{表1 教学评价体系}
\begin{tabular}{p{3.5cm}p{2cm}p{8cm}}
\toprule
\textbf{评价维度} & \textbf{权重} & \textbf{评价要点} \\
\midrule
AI交互日志 & 20\% & 提问的精准性（问题清晰具体）、逻辑性（问题链条完整）、批判性（识别AI不足并修正） \\
\midrule
代码调试记录 & 20\% & 代码理解（准确解释功能原理）、错误修正（独立发现修正错误）、创新改进（优化适应具体情况） \\
\midrule
选址方案报告 & 30\% & 科学性（方法正确、数据规范、权重有据）、创新性（独特视角或方法改进）、图件规范性（要素完整、图例清晰） \\
\midrule
GIS专业报告 & 10\% & 报告规范性（结构完整、写作严谨）、地图制图水平、分析深度 \\
\midrule
反思报告 & 20\% & 反思深度（深入分析AI优势局限）、实践关联（与具体案例结合）、未来展望（建设性建议） \\
\bottomrule
\end{tabular}
\end{table}

评价实施要点：（1）AI交互日志采用档案袋评价，学生全程记录，教师定期检查反馈；（2）代码调试记录重点考查理解深度而非运行成功率；（3）选址方案报告需附GIS操作截图与权重确定论证；（4）反思报告采用答辩形式，学生需现场回答教师提问。

\newpage

\section{六、教学创新点}

\subsection{（一）从"操作导向"到"思维导向"的范式转变}

传统GIS实践教学以软件操作技能为中心，学生学习大量菜单命令但缺乏对空间分析原理的深入理解\citep{汤国安2019}。本案例通过AI承担代码编写、数据处理等技术性任务，将学生从重复性操作中解放，将更多精力投入地理问题的分析与建模。

从认知负荷理论视角，GIS空间分析涉及空间信息感知、地理概念理解、分析流程记忆、软件操作执行等同时处理的认知资源。通过AI承担操作层面的认知负荷，学生认知资源可更多分配到高阶思维任务（问题分析、模式识别、决策判断），提升学习效果与思维深度。

\subsection{（二）"双师协同"教学模式的重构}

本案例构建教师与AI协同的"双师"教学模式：

教师角色重构——从"知识传授者"到"教学设计者"：教师负责设计真实问题情境、规划AI融入节点、制定评价标准，在学生遇到困难时提供启发性提示而非直接答案，引导学生对AI输出进行批判性评估。

AI角色重构——从"辅助工具"到"技术协作者"：AI在"反事实分析"环节快速生成多种情景方案，让学生体会"同一问题、不同策略"的权衡关系；在"反思报告"环节引导学生思考技术应用的伦理边界，培养数字时代的负责任研究意识。

学生角色重构——从"被动接受"到"主动建构"："问题驱动"的AI交互方式比"指令跟随"更能促进深度学习，学生在与AI互动中学会评估输出质量，形成批判性使用AI的意识和能力。

这种人机协同模式形成"人—机—人"互动闭环：学生通过与AI互动发现问题、提出质疑；教师通过观察AI交互了解学习状态、调整教学策略；学生从教师反馈中深化对AI局限性的认识。

\subsection{（三）人机协同素养培育：服务科研能力发展}

本案例注重培养学生的人机协同素养，为未来从事地理信息相关科研工作奠定基础。

培养AI辅助科研能力：学生在本案例中掌握的人机协同工作模式，可直接迁移到未来科研工作中。AI可辅助文献检索、数据处理、代码编写、论文写作等多个环节，显著提升科研效率。

形成批判性使用AI的意识：科研工作中使用AI工具必须保持批判性态度，本案例通过反思报告环节培养学生识别AI局限性的能力，包括数据偏见、算法黑箱、事实性错误等问题。

为GeoAI研究奠定基础：GeoAI是地理信息科学的前沿方向\citep{gai2023}，本案例让学生提前接触AI与GIS的融合应用，为其未来从事GeoAI相关研究奠定基础。

\newpage

\section{七、结语}

本文针对高校地理科学专业GIS实践课的教学困境，设计了基于生成式AI的城市物流园区选址分析教学案例，构建了"问题定义—数据获取—空间建模—成果表达"四阶段人机协同教学模式。

主要贡献包括：（1）提出从"操作导向"向"思维导向"的教学范式转变，将AI工具深度融入GIS实践教学全过程；（2）设计了具体可操作的AI辅助教学流程，提供Prompt示例与任务卡；（3）建立了多元化教学评价体系，突出过程性评价与反思性评价；（4）强调人机协同素养培育，培养学生批判性使用AI工具的能力。

实施建议：本案例对教学条件有一定要求，包括可访问互联网的计算机教室、教师需具备AI工具使用经验、建议安排32-48课时等。未来随着大语言模型与GIS系统的深度集成\citep{gai2023}，学生可能能够通过自然语言直接与GIS系统交互，进一步降低技术门槛。AI辅助的虚拟野外考察、智能地理数据挖掘、个性化GIS学习路径规划等应用场景也值得深入探索。

我们相信，生成式AI不是替代地理研究者的"威胁"，而是拓展地理思维的"外接大脑"。在教育数字化转型背景下，高校地理科学专业应积极拥抱AI技术，培养适应新时代需求的具有数字素养的复合型地理专业人才。

\newpage

\section{参考文献}

\begin{thebibliography}{99}
\addtolength{\itemsep}{-0.5\baselineskip}

bibitem{Jiang2023GeoAI}
Jiang Y, Li Y, Yang C, et al. A comprehensive survey on GeoAI: Foundations, applications, and future trends[J]. International Journal of Geographical Information Science, 2023, 37(8): 1847-1878.

bibitem{杜清运2017}
杜清运， 任福. GIS空间分析的教学内容与方法探索[J]. 地理信息世界， 2017, 24(3): 1-6.

bibitem{Dwivedi2023}
Dwivedi Y K, Kshetri N, Hughes L, et al. ``So what if ChatGPT wrote it?'' Multidisciplinary perspectives on opportunities, challenges and implications of generative conversational AI for research, practice and policy[J]. International Journal of Information Management, 2023, 71: 102642.

bibitem{Zhang2023AI}
Zhang T, Wang D, Yang J. AI-enhanced GIS education: Opportunities and challenges[J]. Journal of Geography in Higher Education, 2023, 47(2): 245-262.

bibitem{Wang2020}
Wang S, Zhang L, Ye Q. The effectiveness of AI-powered language models in GIS education: An empirical study[J]. Computers \& Education, 2020, 145: 103741.

bibitem{Zhao2021}
赵刚， 陈跃， 王吴. 人工智能赋能地理实践教学的模式与路径[J]. 地理教学， 2021(18): 45-49.

bibitem{Luo2022}
罗华， 李文均， 彭岳. 人机协同视域下的地理智慧教育： 理论框架与实践路径[J]. 中国电化教育， 2022(8): 78-85.

bibitem{教育部2022}
教育部. 义务教育地理课程标准（2022年版）[S]. 北京： 北京师范大学出版社， 2022.

bibitem{王民2018}
王民， 蔚东英， 田柳. 中学地理实验教学研究[M]. 北京： 北京师范大学出版社， 2018.

bibitem{李家清2015}
李家清， 常珊珊. 核心素养时代的地理课程与教学改革[J]. 地理教学， 2015(24): 4-8.

bibitem{汤国安2019}
汤国安， 杨昕， 张海平. 地理信息系统教程（第二版）[M]. 北京： 高等教育出版社， 2019.

bibitem{Klosterman2015}
Klosterman R E. Planning support systems and smart cities[M]. Cambridge: Lincoln Institute of Land Policy, 2015.

bibitem{gai2023}
Goodchild M F, Guo H, Maguire M, et al. GeoAI: A perspective on the integration of artificial intelligence and geospatial science[J]. Annals of GIS, 2023, 29(1): 1-7.

bibitem{Zhu2020}
Zhu D, Liu Q, Shen J. The application of GIS in high school geography teaching in China: Current status and future directions[J]. International Research in Geographical and Environmental Education, 2020, 29(4): 345-360.

bibitem{邬伦2014}
邬伦， 刘瑜， 张晶， 等. 地理信息系统——原理、方法和应用[M]. 北京： 科学出版社， 2014.

bibitem{袁孝亭2018}
袁孝亭. 地理学思想与方法[M]. 北京： 高等教育出版社， 2018.

bibitem{中共中央国务院2019}
中共中央，国务院。 中国教育现代化2035[N]. 人民日报， 2019-02-24(001).

bibitem{教育部2020}
教育部. 普通高中地理课程标准（2017年版2020年修订）[S]. 北京： 人民教育出版社， 2020.

\end{thebibliography}

\end{document}